\worldchapter{Creating a Character}{create}
\label{ch:create}

\begin{multicols}{2}

To create a character, you select character templates that reflect their life stages, talents and interests (see \cref{ch:appendix-templates}). Each template can alter the character's attributes and skills, as well as bringing knowledge and shadows.

\wbif{tirakan}{
Additionally, character templates can unlock game mechanics, such as performing the actions of a priest or learning magic spells.
}{}

\section{Reputation}

A character's reputation reflects their level of fame. Characters earn a certain amount of reputation for each adventure they complete. New characters usually start with 20 reputation points. However, the game master can determine this.

Reputation is used to add character templates. Each template costs a certain number of reputation points.

Character templates can have negative reputation costs. In this case, the player receives the points when they select the template. This applies to the \textit{Drunkard} template, for example.

\wbif{phasesix}{
\section{Eras}

Before the campaign or adventure begins, the game master decides which era and extensions to use. This determines which character templates, weapons, armour and items can be used, as well as whether magic, body modifications or priests' actions are possible.

The possible eras are:

\begin{itemize}
    \item Classical antiquity \faIcon{eye}
    \item The Middle Ages, Vikings and Crusades \faIcon{chess-rook}
    \item The Victorian era and the Wild West \faIcon{umbrella}
    \item Imperialism and World Wars \faIcon{flag}
    \item The Cold War and the 1980s \faIcon{plane}
    \item Modern times \faIcon{microchip}
%    \item Near future
    \item Science fiction \faIcon{space-shuttle}
\end{itemize}

Optional extensions include:

\begin{itemize}
    \item Magic \faIcon{magic}
    \item Horror \faIcon{spider}
    \item Pantheon \faIcon{ankh}
    \item Body modifications \faIcon{syringe}
%    \item Vehicles and drones
\end{itemize}

The extensions and eras are listed for all templates, weapons and items, and equipment. 

In the appendix, the symbols used in the list above are used for the character templates, items, weapons, etc.
}{}

\section{Selecting templates}

A character template represents a specific stage in a character's life. Each template is assigned to one of the following categories: education, occupation, talent, interests, character or environment.

Each template alters a small number of the character's attributes and skills, either positively or negatively, and may bring with it knowledge or shadows. Additionally, templates may contain their own rules, which the character then adopts.

\begin{tcolorbox}
\subsection*{Scholar}
\tcblower
\begin{tabular}{@{}ll@{}}
    \textbf{Reputation:} & 10\\
\end{tabular}

\begin{tabular}{@{}ll@{}}
    \textbf{Education} & +4\\
    \textbf{Nature} & +1\\
    \textbf{History} & +2\\
    \textbf{Communication} & +1\\
\end{tabular}
\end{tcolorbox}

Each template is worth a certain amount of reputation. This is the number of reputation points that must be spent to incorporate the template into the character's career.

The list of templates can be found in \cref{ch:appendix-templates}.

\subsection{Base values}

All of a character's attributes, skills and other values start with a uniform base value. Information from the character templates is then added to these values.

\begin{itemize}
    \item Actions: 2
    \item Minimum roll: 5+
    \item Bonus, destiny, and re-rolls: 0
    \item Attributes: 1
    \item Skills: 0
    \item Innate protection: 0
    \item Maximum wouds: 10
    \wbif{nexus}{}{
    \item Arcana: 0
    \item Spell Points: 0
    }
    \wbif{tirakan}{}{
    \item Maximum Stress: 10
    \item Base stress: 0
    \item Biostrain: 0
    }
\end{itemize}

\subsection{Lineage}

First, during character creation, you choose your character’s lineage. This origin not only describes the culture your character comes from, but also provides a lineage template that grants the typical strengths of a member of your people.

Only one lineage template may be chosen, and it does not cost any Reputation.

The available templates are listed in the \cref{ch:appendix-templates}.

The chosen lineage is noted in the career and the specified modifications are added to the character's values.

\subsection{Additional templates}

You can now select as many additional templates as you wish until you have used up all your reputation. You can combine templates from all categories. This means that your character can have one or more occupations, or none at all.

The modifications specified for each template are added to the character's values. In addition, the knowledge, shadows and other rules of the template are added to the character sheet.

All values can also become negative (see \cref{ch:rolls}).

\subsection{Remaining reputation}

Once the player is satisfied with the template, they can declare the character finished. Any remaining reputation that has not been spent will be added to the character's reputation (see \cref{ch:advancement}). This means that no reputation is lost.

\section{Contacts and languages}

Once the character templates have been finalised, the character's languages and contacts can be determined.

\subsection{Contacts}

Contacts are acquaintances or connections that the character had before the start of the campaign. They are recorded with their names and descriptions, and can be imagined as desired.

The number of contacts a new character can have is determined by the sum of the attributes \textit{Charm} and \textit{Attractiveness}.

Contacts are recorded on the character sheet.

\subsection{Languages}

A new character can learn a certain number of languages based on the sum of their \textit{Education} and \textit{Logic} attributes. These can be any languages. If the sum of these attributes is 0 or less, the character has only a limited command of their native language.

Languages are recorded on the character sheet.

\section{Equipment}

Once the character's stats have been determined using the templates, the character can be equipped with gear. The game master sets a starting capital for the characters for the campaign or adventure.

\wbif{tirakan}{
The starting capital is usually 2,000 Gulden.
}{
The starting capital is usually 2,000 units of the standard currency, for example, euros.
}

This starting capital can be used to purchase equipment such as weapons, armor, and items. For more details, see the \cref{ch:equipment} chapter.

\subsection{Equipment}

\cref{ch:appendix-weapons}, \cref{ch:appendix-riotgear} and \cref{ch:appendix-items} can now be purchased with your starting capital. Any purchased items can be noted on the character sheet with their values, and the price can be deducted from your starting capital.

\subsection{Assets}

Any starting capital not spent on weapons, armor, and similar items becomes the character's assets.

\wbif{nexus}{}}
\section{Spells}

\wbif{tirakan}{}{
If the magic extension is used in the adventure or campaign, the character can also learn spells.
}

Character templates offer \textit{spell points} and allow the character to learn spells of a certain \textit{origin}.

If the character has obtained both through the choice of character templates, they can use the spell points to choose spells that they have mastered.

Spells are acquired in a similar way to templates for points. Spell points are used for this purpose. Each spell has a specific cost for which it can be added to the character sheet (see \cref{ch:appendix-spells}). Only spells of origins that the character has unlocked through character templates can be selected. More details can be found in the chapter \cref{cha:chapter-magic}.
}

\wbif{tirakan}{}{
\section{Body Modifications}

\wbif{nexus}{
Unlike the general population, the NEXUS organisation has access to extraterrestrial body modifications. These are mechanical or electronic elements that characters can integrate into their bodies.
}{
When playing with the \textit{Body Modifications} extension, \cref{ch:appendix-bodymodifications} can be purchased and installed for the starting capital.
}

Body modifications can be purchased at the beginning using the starting capital. The rules for body modifications (see \cref{ch:bodymodifications}) must be taken into account here; for example, sufficient energy must be available to power them.

The process of integration by a doctor is not necessary when creating a character; body modifications can simply be noted on the character sheet.
}
\end{multicols}
